% Riferimenti: il nome e il numero sono un unico collegamento cliccabile.
\newcommand{\figref}[1]{\hyperref[#1]{figura~\ref*{#1}}}
\newcommand{\tabref}[1]{\hyperref[#1]{tabella~\ref*{#1}}}
\newcommand{\secref}[1]{\hyperref[#1]{sezione~\ref*{#1}}}
\newcommand{\codref}[1]{\hyperref[#1]{codice~\ref*{#1}}}
\newcommand{\listoflistings}{\listof{codelisting}{Elenco dei codici sorgente}}

% Capitoli finali senza numero, presenti nell'indice e nei segnalibri PDF.
\newcommand{\capitolononnumerato}[1]{%
	\chapter*{#1}%
	\phantomsection
	\addcontentsline{toc}{chapter}{#1}%
}

% Il pedice G si inserisce esplicitamente solo alla prima occorrenza.
% Termine e pedice fanno parte dello stesso collegamento.
\newcommand{\gloss}[2]{\hyperref[glossario-#1]{#2\textsubscript{G}}}
\newcommand{\rimandoglossario}[2]{\hyperref[glossario-#1]{#2}}
\newcommand{\voceglossario}[3]{%
	\item[\protect\phantomsection\label{glossario-#1}#2] #3%
}
